\documentclass[a4paper,10pt]{report}

% Template for the NVKF program

\usepackage{fontspec,xltxtra,xunicode}
\defaultfontfeatures{Mapping=tex-text}
\setsansfont[Scale=MatchLowercase]{Myriad Pro}
\renewcommand{\familydefault}{\sfdefault}
\renewcommand{\ttdefault}{\sfdefault}
\usepackage{multicol}
\usepackage[a4paper, landscape, margin=0.7cm]{geometry}
\usepackage{graphicx}
\usepackage{flushend}
\usepackage{color}
\usepackage{url}
\usepackage{amsmath}
\definecolor{Blue}{rgb}{0.3,0.3,0.9}
\usepackage{polyglossia}
\usepackage{draftwatermark}
\SetWatermarkText{Concept}
\setdefaultlanguage{dutch}

\parskip=0pt

% space between columns
% needs to be twice the margin in order to fold in three
\setlength{\columnsep}{1.4cm}
% width of separation line between the columns
\setlength{\columnseprule}{0.2pt}

\begin{document}
\pagestyle{empty}
\thispagestyle{empty}

\begin{multicols*}{3}

%% Put the logo left top
\noindent
\begin{tabular}{cc}
    \parbox[b]{0.4\linewidth}{%
\includegraphics[width=3cm]{nvkf_logo}
}
& 
%\hspace{-0.5cm}\parbox[b]{0.6\linewidth}{%
\parbox[b]{0.6\linewidth}{%
\begin{center}
{\bfseries Grenzeloze samenwerking}\\ \vspace{0.15cm} 
{\bfseries XXVIVe NVKF-Conferentie}\\ \vspace{0.15cm}
{27 en 28 maart 2015}\\ \vspace{0.15cm}
{Woudschoten Conferentiecentrum}\\
{Woudenbergseweg 54, Zeist}

\end{center}
}
\end{tabular}

%% End of the logo

% To make the header for a plenair item
\newcommand{\plenairheader}[2]{%
\begin{center}
{\color{Blue}{\textbf{#1}}}\\
{\textbf{#2}}
\end{center}
}

% To make the header for a parallel session 
\newcommand{\parallelheader}[2]{%
\parskip=0pt
\noindent{{\color{Blue}{\textbf{\textit{{#1}}}}}\hfill{\textit{{#2}}}}\vspace{-0.2cm}
}

\newcommand{\parallelitem}[1]{%
\noindent \hfill{\color{Blue}{\textbf{\textsc{{#1}}}}}\hfill\strut
}

\newenvironment{packed_enum}{%
    \begin{enumerate}
    \setlength{\itemsep}{1pt}
    \setlength{\itemindent}{0pt}
    \setlength{\parskip}{0pt}
    \setlength{\parsep}{0pt}
    \setlength{\leftmargin}{0pt}
}{\end{enumerate}}

\plenairheader{Plenair Programma}{Vrijdag 27 maart 2015}

\begin{packed_enum}
    \item[09:30] Inschrijving en koffie
        \vfill
    \item[\textbf{10:00}] \textbf{Opening}\\\textit{Richard Kamman, voorzitter NVKF}
        \vfill
    \item[10:05] Grenzeloosheid: zouden we dat willen?\\\textit{Brigit van Ooijen, ziekenhuisapotheker\\Atrium-Orbis MC}
        \vfill
    \item[10:50] Koffiepauze
        \vfill
    \item[{\color{Blue}{\textbf{11:20}}}] {\color{Blue}{\textbf{Bijeenkomsten van de kringen (1)}}}
        \vfill
    \item[12:45] Lunch\\
        {\small Bijeenkomst klinisch fysici in opleiding \hfill\textit{(Plenaire zaal)}}\\
        {\small Bijeenkomst Nederlandse Vereniging voor Klinische Informatica
        in oprichting \hfill\textit{(Zaal 25)}\strut}
        \vfill
    \item[{\color{Blue}{\textbf{14:00}}}] {\color{Blue}{\textbf{Bijeenkomsten van de kringen (2)}}}
        \vfill
    \item[15:30] Theepauze en postersessie 
        \vfill
    \item[\textbf{16:30}] {\textbf{Jonge Onderzoekersprijs (JOP)}}\\\textit{Voorzitter: Maurice Janssen}
    \item[16:30] Point Spread Function modelling improves small lesion detectability in PET/CT\\\textit{Martine Lagerweij}
    \item[16:50] Incorporation of MR Dixon red bone marrow imaging into treatment planning for gynecologic patients treated with chemoradiotherapy\\\textit{Petra Kroon}
    \item[17:10] Geautomatiseerde kwaliteitscontrole van ultrageluidsapparatuur\\\textit{Pepijn van Horsen}
        \vfill
    \item[17:30] Presentatie nieuwe Klifio/KF/KFprof\\Opleidingsinitiatief van het jaar
        \vfill
    \item[\textbf{17:50}] {\textbf{Borreltijd}}
        \vfill
    \item[18:30] Diner
        \vfill
    \item[\textbf{20:30}] {\textbf{Van lab naar veld naar huis: overal meten
        met Rolf Hut}}\\\textit{Rolf ``McGyver'' Hut}
\end{packed_enum}

\columnbreak

\plenairheader{Plenair Programma}{Zaterdag 28 maart 2015}

\begin{packed_enum}

\item[{\color{Blue} {\textbf{09:00}}}]{\color{Blue}{\textbf{Bijeenkomsten van de kringen (3)}}}
        \vfill
\item[10:30] Koffiepauze
        \vfill
\item[\textbf{11:00}] {\textbf{Plenaire sessie}}\\\textit{Voorzitter: Maurice Janssen}
\item[11:00] Radiologist meets physicist: Can we see what we hear?\\\textit{Berit Verbist, LUMC}
\item[11:30] De chemie van samenwerken\\\textit{Otto Bekers, MUMC}
\item[12:00] Physicist meets radiologist: \ldots\\\textit{Andrew Webb, LUMC}
        \vfill
\item[12:30] Sluiting
\item[12:40] Lunch

\end{packed_enum}

\hrule \vspace{3mm}
\vfill
\parallelitem{Parallelsessies (1) vrijdagochtend 27 maart}\\

\parallelheader{Klinische Audiologie}{Zaal 24}
\begin{packed_enum}
\item[\textbf{11:20}] {\textbf{Overzicht wetenschappelijk audiologisch
    onderzoek in de UMC's (1)}}\\\textit{Voorzitters: Bas Franck \& Erwin Baas}
\item[11:20] Onderzoekslijnen in het UMC Utrecht\\\textit{Bert van Zanten}
\item[11:50] Onderzoekslijnen in het LUMC\\\textit{Jeroen Briaire}
\item[12:20] Onderzoekslijnen in het Erasmus MC\\\textit{Andre Goedegebure}
\end{packed_enum}

\vfill
\parallelheader{Algemene Klinische Fysica}{Kapelzaal}
\begin{packed_enum}
\item[\textbf{11:20}]\textbf{Grenzeloze innovatie}\\\textit{Voorzitter: Homme-Auke Kooistra}
\item[11:20] Duo-presentatie Storz/St\"opler en Olympus: 3D endo\-sco\-pie-ontwikkelingen
\item[11:38] Stand van zaken Stille Alarmeringen\\\textit{Carola van Pul}
\item[11:56] Orthopedische koeling\\\textit{Ad Donkerlo}
\item[12:04] RFID door KPN
\item[12:22] Thermal Imaging: (re)discovering the potentials for diag\-nostics and treatment evaluation in the clinical procedures\\\textit{Ruud Verdaasdonk}
\end{packed_enum}
\columnbreak

\parallelheader{Beeldvormende Technieken}{Atrium (A7/A8)}
\begin{packed_enum}
\item[\textbf{11:20}]\textbf{Beeldreconstructiemethoden}\\\textit{Voorzitter: Jan Habraken}
\item[11:20] Introductie, basisprincipe van beeldreconstructiemethoden\\\textit{Antoon Willemsen, UMCG}
\item[11:50] CT beeldreconstructiemethoden\\\textit{Koos Geleijns, LUMC}
\item[12:20] SPECT beeldreconstructiemethoden\\\textit{Hugo de Jong, UMCU}
\item[12:20] PET beeldreconstructiemethoden\\\textit{Johan Nuyts, KUL, Leuven}
\end{packed_enum}

\parallelheader{Klinische Radiotherapie}{Plenaire zaal}
\begin{packed_enum}
\item[\textbf{11:20}]{\textbf{RT Sessie (1)}}\\\textit{Voorzitters: TBA}
\item[11:20] Hypofractionation\\\textit{Jean Philippe Pignol}
\item[11:50] MRI Linac\\\textit{Linda Kerkmeijer}
\item[12:20] Proton Therapy\\\textit{Joachim Widder}
\end{packed_enum}

\parallelheader{Klinische Informatica}{Zaal 25}
\begin{packed_enum}
\item[\textbf{11:20}] \textbf{Handige software voor de klinisch fysicus}\\\textit{Voorzitter: Arjen Becht}
\item[11:20] Digitaal Kernenergiewetdossier succesvol in gebruik\\\textit{Chris Peters, JBZ}
\item[11:50] Update WAD software: kwaliteitsmetingen radiologische apparatuur\\\textit{Joost Kuijer, Vumc}
\item[12:10] Nut en noodzaak voor een dosisregistratiesystemen (inleiding, \textit{Arjen Becht, Gelre Ziekenhuis}):\\
    \vspace{-0.6cm} % bah, vies
    \begin{packed_enum}
    \item[1] Gebruikservaringen commercieel dosisregistratie software pakket \\\textit{Cécile Jeukens, MUMC}
    \item[2] Gebruikservaringen WAD sw dosisregistratie software module \\\textit{Ralph Berendsen, Atrium}
    \item[3] Technische implementatie aspecten en demo REM module in WAD sw  \\\textit{Anne Talsma, MZH}
    \end{packed_enum}
\end{packed_enum}

\columnbreak

\parallelitem{Parallelsessies (2) vrijdagmiddag 27 maart}\\

\parallelheader{Algemene Klinische Fysica}{Kapelzaal}
\begin{packed_enum}
\item[\textbf{14:00}] \textbf{Grenzeloos snijden, coaguleren en ableren}\\\textit{Voorzitter: Job Gutteling}
\end{packed_enum}
\noindent Showcase cauterisatie en andere innovatieve snij-, coagulatie- en
ablatietechnieken. De deelnemers hebben ieder 15 minuten spreektijd. 
\begin{packed_enum}
\item[1] HybridKnife voor gastro-enterologische/urologische procedures\\\textit{Paul Meijers, Erbe}
\item[2] Thunderbeat: (advanced) bipolair als ultrasone techniek\\\textit{Manocher Gholghesaei, Olympus}
\item[3] Plasmajet en AirSeal\\\textit{Don Entjes, Medical Dynamics}
\item[4] Vessel Sealing (Marseal) en bi-polair pincet met irrigatie voor EHT\\\textit{Jos Walvius, KLS Martin}
\item[5] Visie op de toekomst van de Energy Devices\\\textit{Jelle van Riezen, Covidien}
\end{packed_enum}

\vfill 

\parallelheader{Beeldvormende Technieken}{Atrium (A7/A8)}
\begin{packed_enum}
\item[\textbf{14:00}]\textbf{RIVM dosisenquête en wetenschappelijke presentaties BVT}\\\textit{Voorzitter: Cecilia Hindorf}
\item[14:00] Evaluatie RIVM dosisenquête: terug- en vooruitblik\\\textit{Arjen Becht, Gelre Ziekenhuis}
\item[14:20] Toelichting op RIVM analyse CT dosisgegevens: wat kun je er mee?\\\textit{Arnold Schilham, UMCU}
\item[14:30] Development and validation of a patient-tailored dose regime in coronary CT angiography \\\textit{Jorn van Dalen}
\item[14:42] Imaging performance of Ingenuity TF PET/MR versus Gemini TF PET/CT\\\textit{Maqsood Yaqub}
\item[14:54] De productie van 99mTc met een cyclotron\\\textit{Noortje de Groot}
\item[15:06] Quantitative SPECT/CT: where are we, where are we going to?\\\textit{Eric Visser}
\item[15:18] Radiation dose and image quality of novel mammography techniques\\\textit{Leonie Paulis}
\end{packed_enum}
 
\vfill

\parallelheader{Klinische Audiologie}{Zaal 24}
\begin{packed_enum}
\item[\textbf{14:00}] \textbf{Overzicht wetenschappelijk audiologisch onderzoek in de UMC's (2)}\\\textit{Voorzitters: Erwin Baas \& Bas Franck}
\item[14:00] Onderzoekslijnen in het UMC Maastricht\\\textit{Erwin George} 
\item[14:30] Inbedding UAC VUmc, speerpunten/research, samenwerking op campus en (inter)nationaal\\\textit{Theo Goverts}
\item[15:00] Researchlijn op het gebied van spraakverstaan (van Plomp tot nu)\\\textit{Cas Smits}
\end{packed_enum}

\vfill

\parallelheader{Klinische Radiotherapie}{Plenaire zaal}
\begin{packed_enum}
\item[\textbf{14:00}] \textbf{RT Sessie (2)}\\\textit{Voorzitters: TBA}
\item[14:00] Adaptive IGRT\\\textit{Jan Jacob Sonke}
\item[14:30] Automatic Treatment Planning\\\textit{Sebastiaan Breedveld}
\item[15:00] On-line Treatment Planning\\\textit{Bas Raaymakers}
\end{packed_enum}

\vfill
\parallelheader{Klinische Informatica}{Zaal 25}
\begin{packed_enum}
\item[\textbf{14:00}] \textbf{Werken onder architectuur}\\\textit{Voorzitter: Guido Zonneveld}
\item[14:00] Uitleg over referentie domeinen model I-ziekenhuis? / hoe kan dit helpen bij het creëren van overzicht?\\\textit{Hans Boon, Elisabeth ziekenhuis}
\item[14:30] Werken onder Enterprise Architectuur – ervaringen uit de praktijk en wat brengt het ons (niet) / eventueel ook andere sprekers \\\textit{Rob Malschaert, Canisius Wilhelmina Ziekenhuis} 
\item[15:00] Regionale samenwerking op basis van Best Practices: garantie tot succes? \\\textit{Laurens de Groot, Radboud UMC} 
\end{packed_enum}

\vfill
\hrule\vspace{3mm}

\parallelitem{Parallelsessies (3) zaterdagochtend 28 maart}\\
\vfill
\parallelheader{Klinische Audiologie}{Kapelzaal}
\begin{packed_enum}
\item[\textbf{09:00}] \textbf{Vergadering KKau}\\\textit{Voorzitters: Donn\'e Pans en Alex Hoetink}
\item[09:00] Lopende zaken
\item[09:30] Hoofdbehandelaarschap en face-to-face contacten 
\end{packed_enum}
\columnbreak

\parallelheader{Klinische Radiotherapie \& Beeldvormende Technieken}{Plenair}
\begin{packed_enum}
\item[\textbf{09:00}] \textbf{RT Sessie (3)}\\\textit{Voorzitters: TBA}
\item[09:00] MRI for treatment planning\\\textit{Mariëlle Philippens}
\item[09:30] MRI for tumour characterization\\\textit{Uulke van der Heide}
\item[10:00] PET in motion\\\textit{Jeroen van de Kamer}
\end{packed_enum}

\vfill
\parallelheader{Algemene Klinische Fysica}{Atrium (A7/A8)}
\begin{packed_enum}
\item[\textbf{09:00}] \textbf{AKF normen en waarden}\\\textit{Voorzitter: Michael Lansbergen}
\end{packed_enum}
\textit{Sessie in het teken van het cre\"eren van het beeld van het vakgebied Algemene Klinische Fysica en haar professionele randvoorwaarden.}
\begin{packed_enum}
\item[09:00] Grenzeloze samenwerking in de regioziekenhuizen: Van Slotervaart tot Stadskanaal: standaardisatiewerkzaamheden\\\textit{Lieke Poot}
\item[09:18] Stand van zaken kenniscentrum en richtlijn ontwikkeling Federatie Medisch Specialisten\\\textit{Mariken Zijlmans}
\item[09:36] Kwaliteitsnormen t\@.b\@.v\@. uitoefening klinische fysica\\\textit{Homme-Auke Kooistra} 
\item[09:54] Presentatie Werkgroep WINT: Risicomanagement t\@.b\@.v\@. veilig toepassen medische hulpmiddelen\\\textit{Iris Blonk}
\item[10:12] Kwaliteitsnormen t\@.b\@.v\@. scholing van fysische agentia\\\textit{Erik Gelderblom}
\end{packed_enum}

\vfill
\parallelheader{Commissie Stralingshygiëne}{Zaal 24 en 25}
\begin{packed_enum}
\item[\textbf{09:00}] \textbf{Commissie Stralingshygiëne: vervolgsessie CBRN -- oefening opvang radioactief besmette slachtoffers op de SEH}
\end{packed_enum}
\noindent {\small Op 7 november 2014 vond het NVKF symposium "Stralingsincidenten en -ongevallen
bij patiënten" plaats. Klinisch fysici hebben toen in de praktijk kunnen
oefenen wat een stralingsdeskundige moet doen als zich radioactief besmette
slachtoffers melden op de SEH. Ook is een cheat-sheet voor de
stralingsdeskundige opgesteld. Op veler verzoek worden tijdens Woudschoten
deze oefen\-sessies herhaald. Met ervaren CBRN trainers worden alle stappen van
de cheat-sheet doorgenomen en in realistische setting getraind. De cheat-sheet
kan ter voorbereiding hier worden gedownload via:
\url{http://tinyurl.com/k47zcpw}. Let op: deelname is alleen mogelijk voor
degenen die zich hebben ingeschreven via het betreffende onderdeel op het
congres inschrijfformulier.  Opgeven voor wachtlijst via
\url{straling@nvkf.nl}}.

\end{multicols*}

\end{document}
